% Einheit 1 von 5 – Prozente als Anteile (bank/prozentrechnung/e1.jsonl, zusammenbau v0.3)
%% TODO Zweigzeile Teil 1: was hier gelernt wird (Satz in Schülersprache aus dem Einheitstitel) – Einheit 1
\einheitenkopf[e1]{Einheit 1 von 5 · Prozente als Anteile}
\zweigzeile{ab Kl. 6 (am Gymnasium ab Kl. 5) · P10 oft}
\setcounter{aufgabe}{8}

%% TODO Titel: Ich-kann-Satz fehlt in der Bank (Platzhalter: Kettenname) – Nr. 9
%% TODO Anweisung: ein Satz über der ersten Teilaufgabe fehlt in der Bank – Nr. 9
\begin{aufgabe}{Umwandeln}
%% TODO \rechenplatz{n}: Schrittzahl des Grundfalls fehlt in der Bank (nur bei fester Schreibform, 2.3 b)
\begin{teile}
\teil Wie viel Prozent sind grau?

\streifenfeld[2]{50}{$0\,\%$}{$100\,\%$}
\teil $\frac{37}{100}$ – wie viel Prozent? \leerfeld[\%]
\teil $\frac{11}{20}$ – wie viel Prozent? \leerfeld[\%]
\teil $0{,}64$ – wie viel Prozent? \leerfeld[\%]
\teil Jeder hundertste Besucher bekommt einen Preis – wie viel Prozent? \leerfeld[\%]
\teil $35\,\%$ der Kinder fahren Rad, die anderen gehen zu Fuß – wie viel Prozent gehen zu Fuß? \leerfeld[\%]
\teil $5\,\%$ der Pakete kommen zu spät – welche Aussage passt? Kreuze an. \hfill (P10 2020 OS)\\ \kreuz{$5$ von $10$ Paketen kommen zu spät.}\\ \kreuz{Jedes 5. Paket kommt zu spät.}\\ \kreuz{$5$ von $100$ Paketen kommen zu spät.}
\teil Schulweg der Kinder einer Schule. Genau eine Aussage ist falsch. Kreuze sie an und schreibe sie richtig. \hfill (P10 2023 OS)\\ \kreuz{Aussage 1: Ein Viertel der Kinder fährt mit dem Bus.}\\ \kreuz{Aussage 2: Die Hälfte der Kinder fährt mit dem Rad.}

\sachtabelle{lr}{Schulweg & Kinder}{Bus & 150\\ Rad & 200\\ zu Fuß & 180\\ Auto & 70\\ gesamt & 600}
\teil Strom einer Gemeinde nach Energiequelle – trage den fehlenden Anteil ein. \hfill (P10 2021 OS)

\sachtabelle{lr}{Energiequelle & Anteil}{Wind & 38\,\%\\ Sonne & 24\,\%\\ Biogas & \leerzelle\\ Wasser & 9\,\%\\ Sonstiges & 3\,\%\\ gesamt & 100\,\%}
\end{teile}
\end{aufgabe}

%% TODO Titel: Ich-kann-Satz fehlt in der Bank (Platzhalter: Kettenname) – Nr. 10
%% TODO Anweisung: ein Satz über der ersten Teilaufgabe fehlt in der Bank – Nr. 10
\begin{aufgabe}{Prozentangaben ordnen}
\begin{teile}
\teil $0{,}4$; $\frac{1}{4}$; $35\,\%$; $\frac{3}{10}$ – der Größe nach, kleinste zuerst? \leerfeld $<$ \leerfeld $<$ \leerfeld $<$ \leerfeld
\end{teile}
\end{aufgabe}

%% TODO Titel: Ich-kann-Satz fehlt in der Bank (Platzhalter: Kettenname) – Nr. 11
%% TODO Anweisung: ein Satz über der ersten Teilaufgabe fehlt in der Bank – Nr. 11
\begin{aufgabe}{Umwandeln – Fehler finden, Begründen, Darstellungswechsel, Anwendung}
\begin{teile}
\teil Tom: „Jede fünfundzwanzigste Schraube ist fehlerhaft, das sind $25\,\%$.“ – Fehler? Wie heißt es richtig?
\teil Warum ist $\frac{1}{4}$ dasselbe wie $25\,\%$?
\teil Färbe $30\,\%$ des Streifens.

\streifenleer
\teil Abstimmung über den Wandertag: In der 7a sind $18$ von $25$ Kindern für den Zoo, in der 7b $21$ von $30$. In welcher Klasse ist der Anteil für den Zoo größer?
\end{teile}
\end{aufgabe}
